% =============================================================================
% Chapter 14: Method -- Part III (Extended Validation)
% =============================================================================
\mysection{Note on scope (Part III)}
This part reports a later phase of work, conducted after the five/seven-sample PODFAM validation of
Part II, that revisits several of that validation's design choices under controlled ablation: which
loss function the detector is trained with, whether beam-hardening and ring-artefact correction help
on PODFAM data as they did in Part I, whether the circular sample-boundary assumption can be safely
replaced for non-circular parts, and what a genuinely held-out sample reveals about the detector's
real generalisation. It also reports two analyses not tied to any single ablation: a gradient-based
interpretability study of the trained detector, and a strengthening of the pipeline's external
validation against independently published reference porosity. Every figure and table in this part
was generated directly from the working pipeline during this phase; none of Part I's or Part II's own
figures or numbers were re-run, re-verified, or altered here. Where this part's methodology differs
materially from Part II, most importantly, the train/validation/test sample assignment described in
Section~14.3 below, the difference is stated explicitly rather than silently merged with the Part II
description.

\mychapter{14. Method --- Extended Validation}

\mysection{14.1 Scope and Relationship to Parts I and II}
Three questions motivate this part, each building on a specific claim made without controlled
comparison earlier in the thesis. First, Part II's detector (Section~9.6) used the combined Dice-Focal
loss without comparing it against the alternatives Part I's Section~3.8 lists; Section~14.3 below
reports a controlled four-way comparison on PODFAM data. Second, Part I's preprocessing pipeline
(Section~3.3) includes beam-hardening correction (BHC) and ring-artefact suppression, but Part II's
preprocessing (Section~9.2) omits both; Section~14.4 tests whether adding them to the PODFAM pipeline
changes detector performance. Third, Part II's sample-boundary detection (Section~9.3) fits a single
circular mask per volume, an assumption later found not to hold for two PODFAM samples
(sample\_06/07) with non-circular, part-specific cross-sections; Section~14.5 replaces it with a
shape-agnostic method for those samples specifically. A fourth strand, reported in
Sections~14.8--14.10, is not an ablation of an earlier choice but new analysis: which parts of the
trained network drive its output, and how strong is the thesis's external validation once the
reference porosity source's own experimental grounding is checked.

\mysection{14.2 Sample Set: Geometric Diversity and the Waygate Scanner}
All PODFAM samples (Section~9.1) are Ti-6Al-4V, manufactured by laser powder bed fusion on an Aconity
Two machine and scanned on a Waygate v\textbar tome\textbar x M~240 XCT system. The seven-sample set
spans a deliberately varied range of cross-sectional geometry, from the circular cross-sections of
sample\_01--05 to the hexagonal and other non-circular cross-sections of sample\_06 and sample\_07,
included specifically to test the pipeline beyond a single clean cylindrical benchmark. This geometric
diversity is the direct motivation for Section~14.5: a boundary-detection method that assumes a
circular cross-section, adequate for sample\_01--05, systematically misdetects the true material
boundary on sample\_06/07, either clipping real material at the corners of the hexagonal section or
extending the mask past the true boundary into background, which is then misread as defect.

\mysection{14.3 Loss-Function Ablation}
Four loss-function configurations, binary cross-entropy (BCE), Dice, Focal, and combined Dice-Focal
(the formulation used throughout Part II), are trained and compared under an otherwise identical
configuration, isolating loss function as the only varying factor. \texttt{config.py}'s
\param{TRAINING\_SAMPLES\_OVERRIDE} is pinned to the same four-sample list,
\texttt{["sample\_02", "sample\_04", "sample\_03", "sample\_01"]}, most-voids-first, for every run.
\texttt{pipeline.py}'s \texttt{split\_volumes()} assigns roles by rank position rather than by patch,
the first \param{n\_train} names become training samples, the next \param{n\_val} become the single
validation sample, and the remainder become the single held-out test sample, so with four samples and
the pipeline's default 20\%/10\% val/test split this resolves to \textbf{sample\_02 and sample\_04
pooled for training, sample\_03 for validation, and sample\_01 held out entirely for testing}. This is
a whole-sample split, not a patch-level split within a shared pool: sample\_01, the same sample used
as the held-out test role throughout Part II (Section~9.4), never contributes any patch to training
or validation in any of the four runs. sample\_03, excluded from training in Part II as near-empty
(Section~9.4), is used here in the validation role, since a near-empty sample is unsuitable for
training but still usable to monitor convergence. All four runs otherwise use the training
configuration of Section~9.6 unchanged: $256\times256$ patches with 50\% overlap, stratified 1:3
foreground-to-background sampling, Adam optimisation at learning rate $4\times10^{-5}$, up to 50
epochs with early stopping (patience 10 epochs), and best-checkpoint selection by validation Dice.
Every run's total epoch count equals its best epoch plus exactly 10, confirming early stopping, rather
than the 50-epoch cap, was the actual stopping condition in all four cases.

\mysection{14.4 Preprocessing Ablation}
The best-performing configuration from Section~14.3 is retrained with the four-stage Part I
preprocessing pipeline (percentile normalisation, BHC, ring-artefact suppression, NLM denoising,
Section~3.3) applied to the PODFAM volumes in place of Part II's narrower three-stage pipeline
(Section~9.2), using the identical sample split described above. This isolates the marginal effect of
BHC and ring suppression on PODFAM data specifically, since Part I's own parameter sensitivity
analysis (Section~3.3.6) was conducted only on the NIST CoCr dataset and its conclusions were never
previously tested on PODFAM's different scanner and material.

\mysection{14.5 Shape-Aware Sample-Boundary Detection}
For sample\_06 and sample\_07, the circular-fit boundary method (Section~9.3) is replaced with a
sliding-window boundary scan that locates the sample surface by its local intensity peak along each
scan line, without assuming any particular cross-sectional shape, followed by taking the largest
connected component of the resulting interior region and eroding it by a fixed margin
(\param{erosion\_radius} = 35~px) to exclude partial-volume edge pixels. The erosion margin was
selected by a parameter sweep (25, 30, 35, and 40~px), the false boundary artefact it is designed to
remove shrinks smoothly and predictably as the margin increases across this range, and 35~px was the
value visually verified to fully remove the artefact, described in Section~15.3, without eroding
genuine material at the sample interior. A hybrid fallback is required for a minority of slices near
the top and bottom of a volume, where support-structure content rather than open air surrounds the
part: on these transition slices, the primary fixed-threshold detector can fail to find any background
at all (an apparent mask coverage near 100\%), triggering a per-slice adaptive-Otsu re-estimate for
that slice only. This fallback is not fully closed as an edge case, see Section~17 for the specific
failure mode it can itself produce on very high-porosity content. For sample\_01--05, the original
circular-fit method (Section~9.3) is retained unchanged, since it already performs correctly on these
genuinely circular samples; sample\_05 specifically requires a fixed intensity threshold rather than
adaptive Otsu for its own boundary fit, because its unusually high porosity (Section~9.5,
Table~\ref{tab:podfam_classical}) leaves too few solid-material pixels for Otsu's two-class assumption
to separate correctly, a sample-content-dependent choice rather than a universal preference for one
threshold strategy over the other.

\mysection{14.6 A Documented Negative Result: Canny Edge Detection}
Canny edge detection with morphological closing was investigated as a further alternative to both the
circular-fit and shape-aware boundary methods, restricted to boundary detection only (it was not
carried forward into any trained-model comparison). On sample\_01--05, its results closely reproduce
the production circular-fit numbers (Section~15.3), confirming those results independently rather than
replacing them. On sample\_06/07, it exhibits a specific and confirmed failure mode on a small number
of slices: rather than following the true material boundary, it confidently traces the edge of the
surrounding support scaffold instead, producing a boundary that looks well-formed rather than an
obvious error. Investigation traced this to the underlying image signal itself: at the specific slices
where the failure occurs, the material-to-background intensity contrast has genuinely faded to
approximately the same strength as the scaffold's own edge contrast, so no threshold, smoothing
parameter, or shape heuristic (solidity, area continuity, and brightness were each tried as
tie-breakers) can reliably distinguish the two closed shapes from the image content alone. This is
reported as a negative result specifically because the cause is a property of the data at those depths
rather than a fixable implementation choice; the shape-aware method of Section~14.5 remains the
production method for sample\_06/07 for this reason.

\mysection{14.7 Background-Masking Correction for Detector Inference}
A separate defect, found by inspecting the trained detector's predictions on sample\_07 (an
out-of-distribution inference check, not part of the training or validation roles of
Section~14.3), was that the network produced implausible defect predictions on slices with little or
no metal in frame, near-empty transition and support-structure slices the training data never
specifically taught it to recognise as background. The fix applies the same shape-aware boundary
detector used in Section~14.5 to zero out background both before the trained network sees a slice and
again on its predicted output, rather than relying on the network to infer the sample boundary
implicitly from intensity alone.

\mysection{14.8 Conv-Block Attribution Analysis}
To identify which stages of the trained 2D U-Net (architecture per Section~3.7.1, verified by direct
instantiation to have exactly 31,037,633 trainable parameters) contribute most to the defect decision,
a forward hook is registered on every convolutional block's output (four encoder blocks, the
bottleneck, and four decoder blocks, nine blocks total), and \texttt{retain\_grad()} is called on each
hooked output. A forward and backward pass is run on a representative $256\times256$ crop from
sample\_02 (14.27\% genuine Bernsen-labelled defect content), with the scalar objective
$\sum(\hat{y} \odot y)$, the predicted defect probability map $\hat{y}$ elementwise-multiplied by the
Bernsen ground-truth mask $y$ and summed, backpropagated to populate every hooked block's gradient.
Each block's relative contribution is then computed as the mean absolute value of activation times
gradient, $\overline{|a \cdot \nabla a|}$, at that block, normalised across all nine blocks to sum to
100\%. This is a single representative crop and a single attribution method, gradient$\times$activation
is a simple, interpretable choice but not the only one available, so the result reported in
Section~15.6 is directional evidence about where the network's decision-relevant signal concentrates,
not an exhaustive architectural study. The analysis is run on the checkpoint identified in
Section~14.3 as the strongest of the four loss-function configurations (BCE), verified to be the
correct checkpoint by an exact match between its own recorded validation Dice and best epoch
(0.7275 at epoch 13) and the same figures reported independently for the loss-ablation's BCE run.

\mysection{14.9 External Validation Against Independently Published Porosity}
Otsu, Yen, and Bernsen thresholding are each compared, per NIST CoCr sample, against the porosity
values independently published by Kim et al.~\citep{Kim2017} for the same physical samples, using the
Pearson correlation coefficient $r$ across the five samples in common. This validates the classical
methods against a source external to this pipeline, rather than against another output of the same
pipeline. A further check was made of Kim et al.'s own reference values themselves: their paper
reports cross-validating their XCT-derived porosity against an independent gravimetric (density-based)
method, computed from measured disk mass and cobalt-chrome's nominal density, for several of their
samples, with the two methods agreeing to within 1--4\% absolute and an XCT-side threshold-sensitivity
uncertainty below $\pm 0.12\%$~\citep{Kim2017}. This means the reference this thesis's own $r$ value is
measured against is not itself only another XCT measurement, it carries independent physical grounding
one step back. This gravimetric cross-check in Kim et al.'s own paper was reported for specific samples
(their own sample numbering 2, 4, 5, and 6); it has not been independently confirmed here that every
one of the five NIST samples used in this thesis's correlation individually received that exact check,
as opposed to being reported as part of their broader dataset characterisation, and this caveat is
carried into the results and limitations of this part accordingly.

\mysection{14.10 Physical Voxel-Size Determination}
\texttt{config.py}'s \param{PIXEL\_SIZE\_UM} constant, which converts pixel-based measurements
elsewhere in this pipeline to physical units, was never calibrated to either dataset's actual scan
resolution and remains at its unfilled placeholder value of 1.0 (Section~12); this part reports the
actual physical scale for both datasets as a measurement, without editing that constant in the running
pipeline. For PODFAM, the raw TIFF files for sample\_06 and sample\_07 carry embedded
\texttt{XResolution}/\texttt{YResolution} metadata tags, read directly from the file headers, giving a
consistent voxel size of 12.001~\textmu m/px for both samples. For the NIST CoCr samples, the raw TIFF
files carry no equivalent spatial metadata (only an ImageJ-style intensity min/max in each file's
\texttt{ImageDescription} field), so the voxel size is instead taken from the dataset's own published
documentation~\citep{NIST_CoCr}, which states a nominal voxel size of approximately 2.5~\textmu m for
most samples in the set, with samples 3 and 5 specifically scanned at a substantially finer 0.87~\textmu
m. This roughly threefold difference in resolution for exactly the two samples this thesis already
excludes from training, sample\_03 as near-empty and sample\_05 for the DCT labelling artefact
(Sections~9.4, 9.5), is noted in Section~16 as a plausible, previously unidentified explanation worth
stating explicitly, a finer voxel size implies a smaller physical field of view per slice for the same
pixel dimensions, which could plausibly affect both the amount of defect content captured and the
image's intensity/contrast profile. This is reported here as a hypothesis consistent with the
available evidence, not as a confirmed causal finding; it has not been independently tested against
the two samples' actual field-of-view dimensions or intensity statistics within this work.
